% Federated Learning 流程：中央服务器 + 多 client（手机/医院）
% 本地训练 → 上传梯度 → 聚合 → 下发新模型 → 循环
\documentclass[tikz,border=4pt]{standalone}
\usepackage{ctex}
\usepackage{amsmath}
\usetikzlibrary{calc, arrows.meta, positioning, shapes.geometric,
                decorations.pathreplacing}

\begin{document}
\begin{tikzpicture}[scale=1.0, thick,
  server/.style={draw=red!70!black, fill=red!12, rounded corners=5pt,
                 minimum width=2.2cm, minimum height=0.9cm,
                 font=\small, align=center},
  client/.style={draw=blue!60!black, fill=blue!10, rounded corners=3pt,
                 minimum width=1.8cm, minimum height=0.7cm,
                 font=\footnotesize, align=center},
  stepbox/.style={draw=gray!50, fill=gray!10, rounded corners=3pt,
               minimum width=2.0cm, minimum height=0.55cm,
               font=\tiny, align=center},
  upload/.style={-{Stealth[length=5pt]}, orange!80!black, very thick},
  download/.style={-{Stealth[length=5pt]}, blue!60, very thick},
  localstep/.style={-{Stealth[length=4pt]}, green!50!black, thick},
]

% ─────────────────────────────────────────────────────
% 中央服务器（顶部中央）
% ─────────────────────────────────────────────────────
\node[server] (SVR) at (5.5, 6.5) {中央服务器\\（聚合 FedAvg）};

% ─────────────────────────────────────────────────────
% 四个 Client 节点（手机 × 2 + 医院 × 2）
% ─────────────────────────────────────────────────────
\node[client] (C1) at (1.0, 3.8) {手机 A\\（个人数据）};
\node[client] (C2) at (3.5, 3.8) {手机 B\\（个人数据）};
\node[client] (C3) at (7.5, 3.8) {医院 C\\（病历数据）};
\node[client] (C4) at (10.0, 3.8) {医院 D\\（病历数据）};

% ─────────────────────────────────────────────────────
% 上传梯度箭头（橙色，Client → Server）
% ─────────────────────────────────────────────────────
\draw[upload] (C1.north) to[bend left=20]
  node[midway, left, font=\tiny, orange!80!black] {$\Delta\theta_1$}
  (SVR.south west);
\draw[upload] (C2.north) to[bend left=10]
  node[midway, left, font=\tiny, orange!80!black] {$\Delta\theta_2$}
  (SVR.south);
\draw[upload] (C3.north) to[bend right=10]
  node[midway, right, font=\tiny, orange!80!black] {$\Delta\theta_3$}
  (SVR.south);
\draw[upload] (C4.north) to[bend right=20]
  node[midway, right, font=\tiny, orange!80!black] {$\Delta\theta_4$}
  (SVR.south east);

% ─────────────────────────────────────────────────────
% 下发新模型箭头（蓝色，Server → Client）
% ─────────────────────────────────────────────────────
\draw[download] (SVR.south west) to[bend right=20]
  node[midway, right, font=\tiny, blue!70] {$\theta^{r+1}$}
  (C1.north);
\draw[download] (SVR.south) to[bend right=10]
  node[midway, right, font=\tiny, blue!70] {}
  (C2.north);
\draw[download] (SVR.south) to[bend left=10]  (C3.north);
\draw[download] (SVR.south east) to[bend left=20]
  node[midway, left, font=\tiny, blue!70] {$\theta^{r+1}$}
  (C4.north);

% ─────────────────────────────────────────────────────
% 本地训练（Client 内部小循环，绿色箭头）
% ─────────────────────────────────────────────────────
\foreach \c in {C1, C2, C3, C4}{
  \draw[localstep] (\c.south) to[out=270, in=270, looseness=2.5]
    node[below, font=\tiny, green!50!black] {本地 SGD $E$ 步}
    (\c.south);
}

% ─────────────────────────────────────────────────────
% 聚合公式标注（服务器旁边）
% ─────────────────────────────────────────────────────
\node[font=\tiny, red!60!black, align=left] at (9.5, 6.5)
  {$\theta^{r+1} = \sum_k \dfrac{n_k}{N}\theta_k^{r+1}$\\[3pt]（加权平均聚合）};

% 隐私保护标注
\node[draw=gray!40, fill=gray!5, rounded corners=3pt,
      font=\tiny, align=center] at (5.5, 4.9)
  {数据不离本地\\仅上传梯度 $\Delta\theta_k$};

% ─────────────────────────────────────────────────────
% 流程步骤编号（底部）
% ─────────────────────────────────────────────────────
\node[stepbox] (S1) at (1.0, 2.1) {① 下发全局模型};
\node[stepbox] (S2) at (3.8, 2.1) {② 本地训练 $E$ epoch};
\node[stepbox] (S3) at (6.6, 2.1) {③ 上传更新量};
\node[stepbox] (S4) at (9.4, 2.1) {④ 聚合 → 新全局模型};

\draw[localstep] (S1.east) -- (S2.west);
\draw[localstep] (S2.east) -- (S3.west);
\draw[localstep] (S3.east) -- (S4.west);

% 循环箭头
\draw[localstep] (S4.south) to[out=270, in=270, looseness=1.2]
  node[below, font=\tiny, green!50!black] {重复 $R$ 轮}
  (S1.south);

% ─────────────────────────────────────────────────────
% 标题
% ─────────────────────────────────────────────────────
\node[font=\small\bfseries] at (5.5, 7.5)
  {联邦学习（Federated Learning）流程};

\node[font=\tiny, gray!60] at (5.5, 1.0)
  {通信瓶颈 = 轮次数 $\times$ 模型大小；Non-IID 数据导致客户端漂移（client drift）};

\end{tikzpicture}
\end{document}
